\section{INTRODUÇÃO}

A convivência com animais de estimação consolidou-se como um elemento expressivo na estrutura dos lares brasileiros. De acordo com os dados da Pesquisa Nacional de Saúde realizada pelo Instituto Brasileiro de Geografia e Estatística (IBGE, 2020), 46,1\% dos domicílios no país possuíam pelo menos um cão e 19,3\% contavam com ao menos um gato, totalizando dezenas de milhões de animais integrados à rotina das famílias. Essa estreita relação afetiva elevou os cuidados com a saúde e a integridade desses animais, fazendo com que episódios de fuga e desaparecimento causem forte abalo emocional aos tutores e demandem rápida mobilização comunitária (CFMV, 2015).

Apesar da importância dos animais no cotidiano das famílias, a perda acidental de cães e gatos nas cidades continua sendo um problema recorrente e de difícil solução. A circulação desorientada de animais em vias públicas acarreta riscos sanitários, atropelamentos e sobrecarga de abrigos e protetores voluntários (FAROL, 2020). Essa vulnerabilidade ficou ainda mais evidente durante os eventos climáticos ocorridos no estado do Rio Grande do Sul em 2024, quando milhares de animais foram separados de seus tutores durante os resgates de emergência, revelando a fragilidade dos meios tradicionais de identificação e localização (UOL, 2024).

Historicamente, a busca por animais perdidos sempre dependeu da fixação de cartazes em papel nos postes e do contato direto entre vizinhos. Com a popularização dos smartphones, que se tornaram o principal meio de acesso à internet no Brasil (CNDL \& SPC BRASIL, 2024), a divulgação passou a ocorrer também por meio de redes sociais e aplicativos de mensagens instantâneas. No entanto, essas plataformas não foram projetadas para ocorrências de achados e perdidos: as publicações perdem visibilidade rapidamente nas linhas do tempo, não possuem organização geográfica por mapa e muitas vezes expõem dados pessoais e telefones a riscos de golpes e trotes. No cenário atual, ainda há carência de ferramentas gratuitas, colaborativas e focadas na localização imediata do animal.

Diante desse cenário, este trabalho apresenta o desenvolvimento do sistema WeFIND, um aplicativo móvel colaborativo criado para centralizar e agilizar a localização e divulgação de animais perdidos e encontrados, além de apoiar a adoção responsável. A plataforma organiza as ocorrências sobre um mapa interativo em tempo real, permitindo identificar visualmente animais na vizinhança. Para qualificar as buscas, o sistema padroniza o cadastro com fotos e traços essenciais do animal como espécie, raça, porte e pelagem, oferece um canal de mensagens privativo entre os usuários, automatiza a geração de cartazes com código QR para impressão física e disponibiliza uma carteirinha digital de identificação, o RG Pet.

O objetivo central desta pesquisa é desenvolver uma aplicação móvel colaborativa que auxilie tutores e a comunidade na localização e divulgação de animais perdidos e encontrados. Para alcançar esse propósito, o trabalho estrutura-se no cumprimento de metas específicas que envolvem a realização de uma pesquisa prática com a comunidade para identificar as principais dificuldades na busca por animais, a avaliação técnica de sistemas similares existentes no mercado, o levantamento e a documentação detalhada dos requisitos funcionais e não funcionais, a elaboração da modelagem conceitual em notação UML, o projeto da interface centrado na experiência do usuário, a implementação do sistema utilizando React Native e Supabase e, por fim, a realização de testes funcionais e de usabilidade com voluntários do público-alvo.

A relevância social deste projeto está em fornecer uma ferramenta prática e acessível para diminuir o tempo de busca e apoiar a causa animal. Do ponto de vista técnico e acadêmico, o trabalho justifica-se pela oportunidade de integrar conhecimentos de Engenharia de Software, programação para dispositivos móveis, bancos de dados e design de interface desenvolvidos ao longo do Curso Superior de Tecnologia em Sistemas para Internet do IFSul.

A metodologia adotada combina pesquisa descritiva com desenvolvimento tecnológico baseado no modelo incremental, utilizando o método ágil Kanban por meio da ferramenta Trello para o acompanhamento sistemático de cada etapa.

O restante desta monografia está organizado em capítulos sequenciais. O Capítulo 2 aborda o referencial teórico sobre o problema do desaparecimento de animais e o papel das tecnologias móveis em iniciativas comunitárias. O Capítulo 3 detalha a metodologia da pesquisa empírica com tutores, a análise comparativa de soluções similares e o método de desenvolvimento de software adotado. O Capítulo 4 formaliza a engenharia de requisitos funcionais e não funcionais e a modelagem conceitual e lógica do sistema proposto. O Capítulo 5 apresenta formalmente a plataforma WeFIND, detalhando sua arquitetura em nuvem, tecnologias, identidade visual e implementação das telas. O Capítulo 6 documenta os testes de validação funcionais e de usabilidade. Por fim, o Capítulo 7 traz as considerações finais e as sugestões de trabalhos futuros.
\clearpage
